\documentclass[10pt]{article}
\usepackage[margin=0.75in]{geometry}
\usepackage{xcolor}
\usepackage{tcolorbox}
\usepackage{hyperref}
\usepackage{enumitem}
\usepackage{booktabs}
\usepackage{amssymb}

% Define colors
\definecolor{osaorange}{RGB}{220,115,38}
\definecolor{aiblue}{RGB}{30,144,255}

% Configure hyperref
\hypersetup{
    colorlinks=true,
    linkcolor=blue,
    urlcolor=blue
}

% Compact spacing
\setlength{\parindent}{0pt}
\setlength{\parskip}{4pt}
\setlist{noitemsep,topsep=2pt}

\begin{document}

% Orange header
\begin{tcolorbox}[colback=osaorange,colframe=black,boxrule=2pt,arc=0pt,left=6pt,right=6pt,top=8pt,bottom=8pt]
\centering
\textcolor{white}{\textbf{\large WEEK 8 ASSIGNMENT}}\\
\textcolor{white}{\textbf{H524 Introduction to Biostatistics - Fall 2025}}
\end{tcolorbox}

\vspace{4pt}

\noindent\textbf{Total Points:} 10 \quad|\quad \textbf{Due:} Sunday 11:59 PM

\vspace{6pt}

\section*{Part A: Multiple Choice Questions (9 points)}

Answer the following 9 questions on contingency tables and chi-square tests. Each question is worth 1 point.

\subsection*{Question 1: Contingency Table Basics (1 point)}
A contingency table displays the relationship between:
\begin{enumerate}[label=\Alph*., leftmargin=20pt]
\item One categorical variable
\item Two continuous variables
\item Two categorical variables
\item One categorical and one continuous variable
\end{enumerate}

\subsection*{Question 2: Expected Frequencies (1 point)}
In a chi-square test, the expected frequency for a cell is calculated as:
\begin{enumerate}[label=\Alph*., leftmargin=20pt]
\item (Row total × Column total) / Grand total
\item (Row total + Column total) / 2
\item Grand total / Number of cells
\item Observed frequency / Total sample size
\end{enumerate}

\subsection*{Question 3: Chi-Square Test Assumptions (1 point)}
Which of the following is a requirement for conducting a chi-square test?
\begin{enumerate}[label=\Alph*., leftmargin=20pt]
\item All expected cell frequencies must be at least 10
\item All expected cell frequencies should be at least 5
\item The sample must be normally distributed
\item Variables must be continuous
\end{enumerate}

\subsection*{Question 4: Degrees of Freedom (1 point)}
For a 3×4 contingency table (3 rows, 4 columns), what are the degrees of freedom for the chi-square test?
\begin{enumerate}[label=\Alph*., leftmargin=20pt]
\item 7
\item 6
\item 12
\item 11
\end{enumerate}

\subsection*{Question 5: Interpreting Chi-Square Results (1 point)}
A chi-square test yields p = 0.032. At $\alpha = 0.05$, what can we conclude?
\begin{enumerate}[label=\Alph*., leftmargin=20pt]
\item There is no association between the variables
\item There is a significant association between the variables
\item The variables are normally distributed
\item The test assumptions are violated
\end{enumerate}

\subsection*{Question 6: Relative Risk (1 point)}
In a study of disease and exposure, the relative risk (RR) is:
\begin{enumerate}[label=\Alph*., leftmargin=20pt]
\item The ratio of the odds of disease in exposed vs. unexposed groups
\item The ratio of the risk of disease in exposed vs. unexposed groups
\item The difference in risk between exposed and unexposed groups
\item Always equal to 1 if there is no association
\end{enumerate}

\subsection*{Question 7: Odds Ratio (1 point)}
An odds ratio of 2.5 means:
\begin{enumerate}[label=\Alph*., leftmargin=20pt]
\item The exposed group has 2.5 times the risk of disease compared to unexposed
\item The odds of disease in the exposed group are 2.5 times the odds in the unexposed group
\item There is a 2.5\% difference in disease rates between groups
\item The chi-square statistic is 2.5
\end{enumerate}

\subsection*{Question 8: 2×2 Table Setup (1 point)}
In a case-control study examining smoking and lung cancer, which variable should form the rows of a 2×2 table?
\begin{enumerate}[label=\Alph*., leftmargin=20pt]
\item Smoking status (exposed/unexposed)
\item Lung cancer status (case/control)
\item Either variable can form rows
\item Neither - case-control studies cannot use 2×2 tables
\end{enumerate}

\subsection*{Question 9: Fisher's Exact Test (1 point)}
When should Fisher's exact test be used instead of the chi-square test?
\begin{enumerate}[label=\Alph*., leftmargin=20pt]
\item When sample size is very large (n $>$ 1000)
\item When expected cell frequencies are small (less than 5 in one or more cells)
\item When variables are continuous
\item When there are more than two categories
\end{enumerate}

\newpage

\section*{Part B: Group Project Check-In - Progress Update (1 point)}

\begin{tcolorbox}[colback=aiblue!10,colframe=aiblue,boxrule=1pt,arc=2pt,left=6pt,right=6pt,top=6pt,bottom=6pt]
\textbf{GROUP PROJECT PROGRESS CHECK}
\end{tcolorbox}

\vspace{4pt}

\textbf{Reminder:} This Wednesday (12:00-1:20pm) is \textbf{Office Hours Day} - drop by classroom/office if your group needs help!

\subsection*{Your Assignment}

\textbf{Group Leader submits} a brief progress update:

\textbf{Answer these questions:}

\begin{enumerate}
\item \textbf{Has your group met at least once since Week 7?}
\begin{itemize}[label=$\square$]
\item Yes
\item No
\end{itemize}

\item \textbf{Have you loaded and explored your dataset in R?}
\begin{itemize}[label=$\square$]
\item Yes
\item No
\end{itemize}

\item \textbf{Have you consulted at least one AI tool about your data?}
\begin{itemize}[label=$\square$]
\item Yes (which one(s)? \underline{\hspace{3cm}})
\item No
\end{itemize}

\item \textbf{Any blockers or questions?} (optional)\\[4pt]
\underline{\hspace{\textwidth}}\\[2pt]
\underline{\hspace{\textwidth}}
\end{enumerate}

\vspace{6pt}

\textbf{Note:} This week is Office Hours Day! If you answered "No" to any question above or have blockers, come to Wednesday's class time (12-1:20pm) for help.

\subsection*{Submission Format}

\textbf{Submit as:} Text entry (Canvas quiz-style) OR brief paragraph

\textbf{Length:} 3-5 sentences

\textbf{Grading:} 1 point for completion (just answer honestly!)

\subsection*{Notes}

\begin{itemize}[leftmargin=20pt]
\item \textbf{Only the group leader submits}
\item This is a quick check-in (1 point) - not a full report
\item If you're stuck, \textbf{come to Office Hours Day Wednesday 12-1:20pm}
\begin{itemize}[label=--]
\item Location: CORD 1424 (classroom) 12:00-12:30, then 153 Milam (office) 12:30-1:20
\item No appointment needed, just drop by
\item Bring your laptop and dataset
\end{itemize}
\item Groups are expected to be working on analysis this week
\end{itemize}

\end{document}
